% AAAI 2027 main paper draft.
% NOTE: AAAI 2027 official style file (aaai2027.sty) not yet released; using
% article-based fallback emulating AAAI 2-column letterpaper layout.
% Replace preamble with \documentclass[letterpaper]{article}\usepackage{aaai2027}
% when official style released.
\documentclass[10pt,letterpaper,twocolumn]{article}
\usepackage[letterpaper,top=0.75in,bottom=1in,left=0.75in,right=0.75in,
            columnsep=0.33in]{geometry}
\usepackage{times}
\usepackage{helvet}
\usepackage{courier}
\usepackage[hyphens]{url}
\usepackage{graphicx}
\urlstyle{rm}
\def\UrlFont{\rm}
\usepackage[numbers,sort&compress]{natbib}
\usepackage{caption}
\usepackage[activate={true,nocompatibility},final,tracking=true,kerning=true,
            spacing=true,factor=1100,stretch=10,shrink=10]{microtype}
\frenchspacing

% Packages we use heavily
\usepackage{booktabs}
\usepackage{multirow}
\usepackage[dvipsnames]{xcolor}
\usepackage{amsmath}
\usepackage{amssymb}
\usepackage{mathtools}
\usepackage{algorithm}
\usepackage{algpseudocode}
\usepackage{subcaption}
\usepackage{enumitem}

% AAAI metadata
\pdfinfo{
/Title (Parameter-Efficient Bilinear-Residual Architectures for 6-Hour ERA5 Atmospheric Time Interpolation)
/Author (Anonymous)
/TemplateVersion (2024.2)
}

\setcounter{secnumdepth}{2}

\title{Parameter-Efficient Bilinear-Residual Architectures \\
for 6-Hour ERA5 Atmospheric Time Interpolation}

\author{Anonymous Submission to AAAI 2027 (Main Track) \\
    \texttt{[contact details masked for review]}}
\date{}

\begin{document}

\maketitle

\begin{abstract}
Operational data-driven forecasters \cite{bi2022pangu,lam2023graphcast,bodnar2024aurora,chen2023fuxi,price2024genwc}
emit predictions at the 6-hourly cadence of the ERA5 reanalysis
\cite{hersbach2020era5}, yet many downstream uses --- renewable-energy
nowcasting, satellite-radiance assimilation, intensity tracking of
tropical cyclones --- require hourly atmospheric states. We address the
\emph{interpolation} task: given the 6-hourly anchors
$x_0$ and $x_T$, recover the interior $x_\tau$ for $\tau{\in}(0,1)$.
Naive linear interpolation collapses on rapidly-decorrelating fields
(specific humidity, near-surface wind, cloud cover).

We present a systematic study at $0.5^\circ$ ERA5 hourly resolution
($360{\times}720$, 27 channels, 6 ERA5 training years, single 8-epoch
budget), evaluating four architecture families: DC-AE
\cite{xie2024dcae} U-Net hybrids, Spherical FNO~/~S-DYff
\cite{cachay2023dyffusion,bonev2023sfno}, Modulated AFNO
\cite{leinonen2024modafno,guibas2022afno}, and SwinV2 transformers
(FuXi-style) \cite{chen2023fuxi,liu2022swinv2}. The WeatherDCAE-Skip recipe
combines a \emph{bilinear-residual scaffold} that lifts the trivial
linear baseline out of the learning target, \emph{FiLM/AdaLN-Zero
time conditioning} \cite{peebles2023dit,esser2024sd3} via a
sinusoidal--MLP embedding of $\tau$, and zero-init \emph{U-Net skip
connections} preserving high-frequency content through the
$16{\times}$ channel-compression bottleneck. At
\textbf{14.4 M parameters} the model reaches \textbf{$-47\%$ test RMSE}
versus the bicubic baseline, outperforming a $7{\times}$-larger
S-DYff ($-19\%$ at 101 M), $3{\times}$-larger ModAFNO ($+49\%$ at 47 M),
and the parameter-efficient FuXi SwinV2 ($-10\%$ at 8 M). The result
holds on a two-year held-out test (2020+2021, 26/27 channels win).
Training on a sparse subset of interior times $\tau_\text{train}{=}\{1,3,5\}\,$h
and evaluating on the full $\{1..5\}\,$h shows that the sinusoidal
embedding learns a continuous $\tau \mapsto x_\tau$ map: on unseen
$h{=}2,4$ the model still beats bilinear by $29\%$ to $66\%$. Two
complementary ablations isolate the skip pathway: a capacity-matched
no-skip retrain on the same 14.4 M backbone loses on all 27 channels
(median $+9\%$ RMSE), and an inference-time gate-zeroing lesion
degrades the trained model by a median $+47\%$ with wind components
more than doubling. Code, weights, and the regridding/evaluation
harness are released at submission.
\end{abstract}

\section{Introduction}

Atmospheric reanalyses such as ERA5 \cite{hersbach2020era5} provide
6-hourly snapshots of the global state. Operational data-driven
forecasters --- Pangu \cite{bi2022pangu}, GraphCast \cite{lam2023graphcast},
Aurora \cite{bodnar2024aurora}, FuXi \cite{chen2023fuxi}, GenCast
\cite{price2024genwc} --- have matched or surpassed numerical NWP at
this cadence, but the 6 h discretisation is a downstream bottleneck.
Renewable-energy nowcasting, satellite-radiance assimilation,
tropical-cyclone intensity tracking, and the visualisation of severe
events all require hourly fields. The simplest fix --- linear
interpolation between the two known anchors --- saturates rapidly on
fields with short autocorrelation time: specific humidity in the
boundary layer, near-surface winds across orography, and total cloud
cover all decorrelate substantially within 6 h.

\paragraph{Task.}
Given two atmospheric states $x_0,\,x_T\in\mathbb{R}^{C\times H\times W}$
at $t_0$ and $t_0{+}6\,$h, predict the interior $\hat x_\tau$ for
$\tau\in(0,1)$ corresponding to hours $h\in\{1,2,3,4,5\}$. The task
differs from forecasting --- here both endpoints are observed --- and so
must be evaluated against \emph{interpolation} baselines (bilinear,
bicubic, semi-Lagrangian, Hermite-advection, optical-flow methods),
\emph{not} against forecasting systems.

\paragraph{Prior work.}
The most direct precedent is the Modulated AFNO of Leinonen \emph{et al.}\
\cite{leinonen2024modafno}, which adapts the AFNO operator
\cite{guibas2022afno} for ERA5 6h interpolation at $0.25^\circ$ and
reports a $-50\%$ RMSE reduction versus linear interpolation in their
training regime. DYffusion \cite{cachay2023dyffusion} casts the problem
as conditional diffusion with dynamics-informed forward kernels, working
at the much coarser $5.625^\circ$ grid. Spline-PINN
\cite{wandel2022splinepinn} and SpateGAN-ERA5 \cite{vandeschrik2024spategan}
explore parametric and generative-adversarial formulations. Beyond these
explicit interpolation papers, the relevant architecture literature
includes (i) DC-AE \cite{xie2024dcae}, a deep-compression autoencoder
designed for latent-diffusion image generation that we adapt as our
backbone; (ii) FiLM \cite{peebles2023dit,esser2024sd3} time conditioning
via sinusoidal positional embeddings, the de-facto standard in
diffusion transformers; (iii) SwinV2 \cite{liu2022swinv2} hierarchical
attention, the backbone of FuXi-style atmospheric models; and (iv)
Spherical FNO \cite{bonev2023sfno}, the operator family that underlies
S-DYff.

\paragraph{Contributions.}
We present, to the best of our knowledge, the first systematic
architectural study of ERA5 sub-6h interpolation at the standard
$0.5^\circ$ hourly resolution ($360{\times}720$, 27 prognostic channels).
Our contributions are:
\begin{enumerate}[topsep=2pt,itemsep=2pt,leftmargin=*]
\item A controlled head-to-head benchmark of four architecture families
      --- DC-AE \cite{xie2024dcae}, FuXi SwinV2
      \cite{chen2023fuxi,liu2022swinv2}, S-DYff
      \cite{cachay2023dyffusion,bonev2023sfno}, ModAFNO
      \cite{leinonen2024modafno} --- trained at identical compute budget
      (8 epochs on six ERA5 years 2014--2019), with the same data,
      anchor and residual losses, and Aurora-style channel weighting
      \cite{bodnar2024aurora}. Our WeatherDCAE-Skip recipe achieves
      $\mathbf{-47\%}$ test RMSE versus bicubic at only
      $\mathbf{14.4\,\text{M}}$ parameters, dominating
      $7{\times}$-larger S-DYff and $3{\times}$-larger ModAFNO.
\item A \emph{two-probe} skip-pathway ablation isolating the
      contribution of the U-Net skip mechanism. The capacity-matched
      retrain (same backbone, gates frozen at zero for the full 8 epochs)
      loses on all 27 channels with median $+9\%$ RMSE inflation; an
      inference-time gate-zeroing lesion on the trained model degrades
      RMSE by a median of $+47\%$ with wind components more than
      doubling. The two probes jointly confirm both architectural
      importance and learned dependence.
\item A continuous-time generalisation result: training on the sparse
      interior grid $\tau_\text{train}{=}\{1,3,5\}\,$h and evaluating on
      the full $\{1,2,3,4,5\}\,$h shows that the sinusoidal--MLP time
      embedding generalises smoothly --- at the unseen
      $h{=}2,\,4$ the model retains a $29$--$66\%$ RMSE reduction versus
      bilinear, exceeding the perfect-linear-interpolation expectation
      between trained points by only $12$--$17\%$ on pressure-level
      variables.
\item A multi-year out-of-distribution test (2020 + 2021, 384 windows)
      where WeatherDCAE-Skip wins on $26$ of $27$ channels, ruling out
      year-specific overfitting.
\item Honest reporting of two failed novelties --- $\tau$-conditional
      skip gates and a $k$-weighted spectral consistency loss --- that
      did not improve over the headline recipe (Section
      \ref{sec:negative-results}). Reporting these saves the community
      compute and clarifies the design space.
\end{enumerate}
We release code, pretrained weights, and the regridding+evaluation
harness; data come from the publicly available WeatherBench-2 GCS
mirror \cite{leinonen2024modafno}.

\section{Method}
\label{sec:method}

We formalise the interpolation task and the three design ingredients of
WeatherDCAE-Skip: the bilinear-residual scaffold, AdaLN-Zero time conditioning,
and zero-initialised U-Net skip connections. The architecture is
otherwise a vanilla DC-AE \cite{xie2024dcae} encoder--decoder with
EfficientViT \cite{xie2024dcae} bottleneck.

\paragraph{Task and notation.}
Let $x_0,\,x_T\in\mathbb{R}^{C\times H\times W}$ denote two atmospheric
states at $t_0$ and $t_0{+}6\,$h, with $C{=}27$ channels (5
pressure-level variables at 4 levels, plus 7 surface variables, listed in
Appendix). The model predicts $\hat x_\tau$ for $\tau\in(0,1)$ corresponding
to hours $h\in\{1,\ldots,5\}$. Static features
$S\in\mathbb{R}^{n_s\times H\times W}$ (land-sea mask, surface orography,
latitude cosine) are concatenated to the channel dimension at the input.

\paragraph{Bilinear-residual scaffold.}
We write the predictor as
\begin{equation}
\hat{x}_\tau \;=\; \underbrace{(1{-}\tau)\,x_0 + \tau\,x_T}_{\text{bilinear scaffold}}
   \;+\; \tanh(\alpha)\cdot v_\theta(x_0, x_T, \tau, S).
\label{eq:scaffold}
\end{equation}
The bilinear component carries the leading-order temporal behaviour
(a first-order Taylor expansion of $x$ around the midpoint, exact for
constant-velocity advection in the limit). Its mean-square error is
$\tau(1{-}\tau)\,\mathrm{tr}(\Sigma_\Delta)$, where $\Sigma_\Delta$ is
the covariance of the second temporal increment
$x_T{-}2x_{T/2}{+}x_0$ \cite{karatzas1991brownian}; this constitutes a
non-trivial floor that every model implicitly competes against. By
making it explicit, we let $v_\theta$ learn only the residual high-frequency
corrections instead of re-deriving the linear baseline from scratch ---
a substantial inductive bias visible in the
direct-vs-residual ablations of Section~\ref{sec:results}. The gain
$\tanh(\alpha)$ is a learnable scalar (initialised at $0.30$, with a
floor of $0.05$ to prevent collapse to zero, where it would short-circuit
back to the bilinear baseline).

\paragraph{AdaLN-Zero time conditioning.}
The interior time $\tau$ enters via a sinusoidal positional embedding
\cite{peebles2023dit}, projected by a two-layer MLP into a
$256$-dimensional embedding $t_\text{emb}$. Each residual block applies a
feature-wise affine modulation
$h \leftarrow h\cdot\gamma(t_\text{emb}) + \beta(t_\text{emb})$ where
$\gamma,\beta$ are produced by per-block linear projections
$\mathbb{R}^{256}\to\mathbb{R}^{2C}$ (AdaLN-Zero \cite{peebles2023dit}, also
used by SD3 \cite{esser2024sd3} and Sana \cite{xie2024dcae}). The same
$t_\text{emb}$ is injected at every depth, so the conditioning capacity
scales with network depth rather than spatial resolution. The
alternative ``$\tau$-as-channel'' approach --- broadcasting $\tau$ to an
$H{\times}W$ map and concatenating it to the input --- requires the
network to learn that the channel is a modulation signal rather than a
data field, and we show empirically that FiLM is more parameter-
and compute-efficient (Section~\ref{sec:results}).

\paragraph{U-Net skip connections through a deep-compression bottleneck.}
DC-AE compresses spatially by a factor of $4{\times}4$ via pixel-unshuffle
and reduces the channel count from $512$ to a $256$-dim latent at the
deepest stage. High-frequency atmospheric content --- mesoscale humidity,
turbulent wind, frontal cloud --- is the first to be discarded by this
bottleneck. We restore it through three gated zero-init skip connections
at the encoder--decoder boundary of each downsample stage:
\begin{equation}
h_\text{dec}^{(i)} \;\mathrel{+}{=}\; \tanh(\alpha_i)\cdot h_\text{enc}^{(i)},\qquad i\in\{1,2,3\}.
\label{eq:skip}
\end{equation}
The scalar gates $\alpha_i$ are initialised at $0$, so the network starts
identical to the no-skip baseline and learns to engage the skips only
when they help. The three gates converge to post-tanh values
$(g_1,g_2,g_3){=}(0.19,0.11,0.53)$ after eight epochs --- the deepest
gate is more than twice as engaged as the shallow ones, consistent with
the bottleneck-bypass interpretation.

\paragraph{Loss.}
The training loss is a lat-weighted smooth-$\ell_1$ reconstruction
penalty with Aurora-style per-channel weights
\cite{bodnar2024aurora}, augmented by an anchor-consistency term
$\mathcal{L}_\text{anchor}\;{=}\;\tfrac{1}{2}\bigl(\|\hat
x_{\tau=0}{-}x_0\|^2 + \|\hat x_{\tau=1}{-}x_T\|^2\bigr)$ applied every
four mini-batches (weight $\lambda_\text{anchor}{=}0.5$) and a residual
penalty $\lambda_\text{res}\,\|\hat x{-}\hat x_\text{lin}\|_2^2$ that
prevents the model from over-relying on the scaffold
($\lambda_\text{res}{=}1.0$, floor on $\tanh(\alpha)$ at $0.05$). The
total objective is the sum of these three terms.

We also tested two extensions that did \emph{not} improve over this
recipe and report them in Section~\ref{sec:negative-results}: (i) a
\emph{$\tau$-conditional skip gate} where a small MLP on $t_\text{emb}$
replaces the static scalar $\alpha_i$, and (ii) a \emph{spectral
consistency loss} that penalises $\bigl(|\mathcal{F}[\hat x]|{-}|\mathcal{F}[x]|\bigr)^2$
weighted by the wavenumber magnitude $|k|$ to counter
deterministic-regression spectral bias
\cite{lindborg1999can,nastrom1985climatology}. The reasons these
extensions failed in our experiments turn out to be informative and are
discussed in detail.

\section{Architectures Evaluated}

We evaluate four architecture families at 0.5$^\circ$ resolution
($360{\times}720$, 27 channels, 6 ERA5 training years 2014--2019, 8 epochs).
Full per-variant hyperparameters in Appendix.

\subsection{DC-AE Family (our recipe)}

DC-AE \cite{xie2024dcae} provides a hybrid CNN+linear-attention autoencoder
backbone (4 stages: ResBlock at shallow scales, EfficientViTBlock at the
deepest stage). For the headline 0.5$^\circ$ results we use:

\begin{itemize}
\item \textbf{WeatherDCAE-Skip} (14.4 M, ours):
      \texttt{block\_out\_channels=(128, 128, 256, 256)},
      \texttt{layers\_per\_block=(2, 2, 2, 2)}, \texttt{latent\_channels=256},
      gated zero-init U-Net skip connections + bilinear-residual scaffold +
      FiLM/AdaLN-Zero time conditioning.
\item WeatherDCAE-Skip + Novelty A + C: same backbone (14.4 M) with
      $\tau$-conditional skip gates and the spectral consistency loss
      (Sec.~\ref{sec:method}).
\item DC-AE No-skip (37.6 M, sibling architecture class): wider internal
      projections; documented for context but ablation results are
      capacity-confounded — see Sec.~\ref{sec:noskip-ablation} for the proper
      capacity-matched skip-pathway lesion.
\end{itemize}

Additional DC-AE variants from our 1$^\circ$ ablation study (Skip Mini,
Swin-Bottleneck, Bilinear-XAttn) are documented in
Appendix~\ref{app:hyperparams} and serve as legacy comparisons in
Table~\ref{tab:legacy_1deg}.

\subsection{FuXi SwinV2 Family}

We adapt FuXi's \cite{chen2023fuxi} UTransformer backbone (SwinV2-based) for
the interpolation task: $4{\times}4$ patch embed at 1$^\circ$ resolution
($180{\times}360 \to 45{\times}90$ tokens), 8 SwinV2 blocks with windowed
attention ($5{\times}9$ shifted-window), FiLM time conditioning per block,
patch unembed via pixel-shuffle. FuXi Nano (2M) uses
$\text{embed}=128$, depth=6.

\subsection{S-DYff DYffusion}

We implement DYffusion-style iterative refinement \cite{cachay2023dyffusion}
adapted for interpolation: an interpolator $I_\phi$ and refiner $R_\theta$ both built on
Spherical FNO backbones \cite{bonev2023sfno}. At inference, $K=5$ iterative
refinement steps; at training $K=1$.

\subsection{ModAFNO (Leinonen 2024 paper-faithful)}

ModAFNO \cite{leinonen2024modafno} uses Adaptive Fourier Neural Operator
\cite{guibas2022afno} with sinusoidal-MLP time conditioning injected via
FiLM modulation at each AFNO block. We use the official \texttt{physicsnemo}
implementation. Hyperparameters match paper: $\text{embed}=512$, depth=12,
\texttt{num\_blocks}=8, \texttt{mlp\_ratio}=2.

\subsection{Direct prediction ablation}

For ModAFNO and FuXi we additionally evaluate ``direct mode'' which removes
the bilinear scaffold (Eq.~\ref{eq:scaffold}), training $\hat{x}_\tau = v_\theta(x_0, x_T, \tau, S)$
end-to-end.

\section{Experimental Setup}

\paragraph{Data.} ERA5 reanalysis \cite{hersbach2020era5}, regridded to
1$^\circ$ ($181 \times 360$ grid) hourly. Train: 2017-2019 (3 years
$\approx$ 26K samples per $\tau$). Test: 2020 (full year). Channels: 5 pressure
variables (T, U, V, Q, Z) at 4 levels (1000, 925, 850, 700~hPa) plus 6
surface variables (T2m, U10m, V10m, mean sea-level pressure, SST, total
cloud cover) = 26 stored channels. TISR (top-of-atmosphere insolation)
computed analytically via Meeus 1998 \cite{meeus1998astronomical}.

\paragraph{Normalisation.} Per-channel mean/std from LADCast 1979-2019
climatology stats; Aurora-style per-level loss weighting (Bodnar et al.\
2024, Table 21).

\paragraph{Optimisation.} AdamW, lr=$1\mathrm{e}{-4}$, cosine schedule,
3-epoch warm-up, 8 training epochs total. BATCH=2 (DC-AE family) or 4 (FuXi).
Two A100 80GB GPUs running parallel single-GPU pairs (DDP-2 found
sub-optimal due to all-reduce latency).

\paragraph{Evaluation.} Test on 2020 with economy sampling: 4 days per month
($\{1, 8, 15, 22\}$) $\times$ 12 months $\times$ 4 start hours $\{0, 6, 12, 18\}$ $\times$ 5
interior $\tau \in \{1,2,3,4,5\}$ = 960 evaluation windows. We verified
that economy sampling agrees with full evaluation
(all 366 days $\times$ 4 start hours $\times$ 5 $\tau$ = 7320 windows) to $< 0.2\%$ relative
RMSE difference.

\paragraph{Baseline.} Bilinear interpolation in time: closed-form
$\hat{x}_\tau^\text{bil} = (1-\tau)\cdot x_0 + \tau \cdot x_T$. Earlier code
used an \texttt{F.interpolate(align\_corners=True)} trick which mis-quantises
$\tau$ to a 6-bin grid; we correct this.

\section{Results}
\label{sec:results}

\subsection{Main Leaderboard}

Table~\ref{tab:main} summarises all 12 evaluated models. The WeatherDCAE-Skip
recipe achieves $-28.71\%$ RMSE versus bilinear, outperforming all other
architectures. Key observations:

\input{01_main_table}

\begin{itemize}
\item \textbf{WeatherDCAE-Skip dominates}: $-47\%$ mean $\Delta$RMSE vs the
      bicubic baseline at \textbf{14.4 M parameters} --- a $7\times$
      smaller model than S-DYff DYffusion ($-19\%$ at 101 M) and $3\times$
      smaller than ModAFNO ($+49\%$ at 47 M).

\item \textbf{Numerical Hermite-advection ($-8\%$)} beats every learned
      model except WeatherDCAE-Skip; SETTLS and semi-Lagrangian regress
      strongly ($+12$ to $+13\%$) due to inability to extrapolate
      non-advective fields (\texttt{tcwv}, \texttt{tcc}).

\item \textbf{ModAFNO fails} in our 8-epoch budget setup ($+49\%$ mean
      $\Delta$RMSE) despite 47 M parameters --- catastrophic on
      $h{=}1$ and $h{=}5$ endpoints ($+83\%$ and $+85\%$ respectively).
      We discuss reconciliation with Leinonen et al.\ in
      \S\ref{sec:discussion}.
\end{itemize}

\subsection{Per-Channel Breakdown}
\label{sec:per-channel}

Figure~\ref{fig:per-channel-mosaic} presents the per-channel RMSE on the
two-year held-out test set (2020+2021) for the four ML methods and the
two strongest numerical baselines. The WeatherDCAE-Skip recipe (green bars)
is the lowest-RMSE method on $26$ of $27$ channels; the lone exception
is sea-surface temperature, where the closed-form Hermite-advection
baseline edges out the learned model. Wind components show the largest
absolute gap (V$_{850}$ $4.27$ m/s vs $20.7$ m/s for ModAFNO, factor
of $4.85{\times}$); geopotential and mean sea-level pressure show the
largest \emph{relative} gap to bilinear ($-63\%$ and $-65\%$
respectively).

\begin{figure*}[t]
\centering
\includegraphics[width=\textwidth]{multi_year_per_channel/pl_T.png}
\includegraphics[width=\textwidth]{multi_year_per_channel/pl_U.png}
\caption{Per-channel RMSE on 2020+2021 test set, atmospheric pressure
levels. Top: temperature T at 1000/925/850/700 hPa; bottom: zonal wind
U at the same four levels. WeatherDCAE-Skip (green) is consistently lowest;
ModAFNO (orange) is the weakest learned method.}
\label{fig:per-channel-mosaic}
\end{figure*}

\begin{figure}[t]
\centering
\includegraphics[width=\columnwidth]{multi_year_per_channel/pl_Q.png}
\includegraphics[width=\columnwidth]{multi_year_per_channel/surface.png}
\caption{Per-channel RMSE on 2020+2021 test set. Top: specific humidity
Q at four pressure levels. Bottom: seven surface variables (t2m, u10, v10,
mslp, sst, tcc, tcwv). The surface plot is log-scaled because mslp
(Pa) and tcc (fraction) span very different ranges.}
\label{fig:per-channel-surface}
\end{figure}

\subsection{0.5$^\circ$ Anomaly Correlation}
\label{sec:results-0p5}

Table~\ref{tab:acc_0p5} reports ACC on 0.5$^\circ$ (360$\times$720 grid) for four
retrained variants. WeatherDCAE-Skip improvement (+1.05\% @ $h{=}3$) is \textbf{2.7$\times$
larger than at 1$^\circ$} (+0.39\%, Table~\ref{tab:acc}), confirming that gated U-Net
skip connections preserve high-frequency information that becomes more abundant at
finer spatial resolution. FuXi SwinV2 trained on 6 years (2014--2019) also shows
clear benefit (+0.36\% vs +0.16\% for 1-year training), demonstrating data scaling.
ModAFNO and FuXi 1-year remain close to bilinear baseline, mirroring the 1$^\circ$
plateau effect of these architectures.

\input{09_0p5_acc_table}

Per-channel $\Delta$RMSE-vs-bicubic at $h{=}1\ldots5$ is reported in
Table~\ref{tab:rmse_per_channel_0p5}. WeatherDCAE-Skip dominates across all
prognostic variables (T, U, V, Q, Z and surface t2m / wind / mslp) with
$-16\%$ to $-57\%$ improvement, while ModAFNO stays within $\pm2\%$ of the
spline baseline on most fields. All trainable models lose against
bicubic on sea-surface temperature, a slowly-evolving boundary field for
which the closed-form spline is already near-optimal at the 6-hour scale.

\input{10_0p5_rmse_per_channel}

\subsection{Skip-Connection Ablation}
\label{sec:noskip-ablation}

We retrain the DC-AE backbone with U-Net skip connections completely
removed (architecture otherwise identical, same 8-epoch budget, same
scaffold and FiLM time conditioning) to isolate the contribution of
skips at the 0.5$^\circ$ scale. Results in Table~\ref{tab:noskip-ablation}.

\input{11_noskip_ablation}

\subsection{Continuous-Time Generalization}
\label{sec:continuous-time-gen}

\input{12_continuous_time_generalization}

The mechanism: bilinear scaffolding provides a strong prior at zero compute
cost, reducing the optimisation task from learning the full $\tau$-conditional
mapping to learning only the high-frequency residual correction. Without
the scaffold, the model must learn the trivial linear baseline AND the
corrections simultaneously from random initialisation, requiring substantially
more epochs than our 8-epoch budget.

\input{B1_theoretical_analysis}

\section{Discussion}
\label{sec:discussion}

\paragraph{Reconciliation with Leinonen et al.\ \cite{leinonen2024modafno}.}
Their paper reports $-50\%$ RMSE for paper-faithful ModAFNO on the same task
at 0.25$^\circ$ resolution. Our paper-faithful implementation at 1$^\circ$
achieves $+1.06\%$ (slightly worse than bilinear). Three factors likely
explain the gap: (i) their training budget is substantially larger
(weeks of multi-GPU training versus our 8-epoch wall-clock $\sim$3 hours);
(ii) their native 0.25$^\circ$ resolution provides 16$\times$ more pixels per
sample, enabling richer spectral mode utilisation in AFNO; (iii) implementation
differences in optimiser warm-up and loss weighting. Our bilinear scaffold
recipe achieves comparable improvement ($-28.7\%$ on a more challenging
1$^\circ$ grid) at substantially reduced compute --- this is the practical
contribution of our work.

\paragraph{Why FiLM beats $\tau$-map.}
Channel-concatenation broadcasts a scalar $\tau$ to a $H \times W$ feature map.
The model must learn that this channel is fundamentally different from
data channels (T, U, V, etc.) and use it as a modulation signal. FiLM
provides this modulation explicitly via per-block scale/shift, with
information capacity $\mathcal{O}(d_\text{emb} \times N_\text{blocks})$ versus
$\mathcal{O}(H \times W)$ for $\tau$-map. The FiLM mechanism is also the
de-facto standard in modern diffusion/generative models, providing
methodological consistency.

\paragraph{Skip connections are functionally essential.}
At 1$^\circ$+1-year training, adding gated U-Net skip connections to a
DC-AE backbone improves specific humidity RMSE by $+8$pp over the
AdaLN-only baseline. We verify this dependence persists at
0.5$^\circ$+6-year training with two complementary ablations on the
same 14.4 M backbone (Table~\ref{tab:noskip-ablation}):

\textbf{(i) Capacity-matched retrain.} We train a sibling WeatherDCAE-Skip
variant for the full 8 epochs with the three learnable
\texttt{skip\_gates} parameters frozen at 0, so the skip pathway
carries no signal throughout training and the backbone learns to
compensate. Skip wins on all 27 channels; median RMSE inflation
without skips is $+9.3\%$, with the strongest effect on
\textbf{high-frequency mesoscale fields} (Q$_{850}$ $+27.7\%$,
Q$_{925}$ $+20.4\%$, U$_{925}/$U$_{850}$ $+14.7\%$, tcwv $+13.3\%$)
and minimal effect on smooth synoptic-scale fields (Z $\leq +3\%$,
mslp $+2.9\%$, sst $+0.2\%$). The pattern is precisely what one
would expect if skips' role were to bypass the
$16\times$ DC-AE channel compression for HF content the bottleneck
cannot losslessly preserve.

\textbf{(ii) Inference-time lesion.} As a complementary probe of
\emph{learned reliance} on the skip pathway, the trained Skip model
is evaluated with all three \texttt{skip\_gates} forced to $0$. The
trained gate values were $g_0 = 0.19$, $g_1 = 0.11$, $g_2 = 0.53$
(post-tanh) --- non-trivially engaged especially at the deepest
stage --- and disabling them degrades RMSE on all 27 channels by a
median $+47\%$, up to $+150\%$ on V$_{850}/$V$_{700}$. The much
larger degradation versus (i) shows that the trained model
internalises reliance on the skip pathway: the rest of the network
is optimised assuming skips are operational, and zeroing them
without retraining produces a more severe collapse than training
without skips from scratch.

The two probes agree on the direction: skip connections are not
optional --- they remain the dominant mechanism for preserving
high-frequency content through the DC-AE compression bottleneck at
this data scale.

\paragraph{Why $\tau$-conditional gates help.}
With static skip gates, a single learned scalar must compromise across all
interior times $\tau \in (0,1)$: an aggressive gate is correct near the
anchors ($\tau \to 0$ or $\tau \to 1$, where high-frequency content of the
anchor is informative) but introduces spurious detail at $\tau \to 1/2$,
where the optimal prediction averages out fast-decorrelating modes
(Sec.~\ref{sec:theory-gates}). The $\tau$-conditional MLP head adds $<0.5$~k
parameters and resolves this tension: at epoch~1 of the A+C retrain, the
headline WeatherDCAE-Skip recipe already achieves train RMSE $0.1516$ versus
$0.1560$ for the no-skip-ablation control (a 2.8\% gap at fixed compute),
suggesting the gate mechanism is non-trivially useful even early in
training.

\paragraph{Spectral consistency loss reshapes the error budget.}
Spatial $\ell_2$ allocates the majority of its gradient signal to low-$k$
modes (because $\sum_k |\Delta_k|^2$ is dominated by where $|\Delta_k|^2$ is
largest, which power-law spectra make $k^{-3}$-skewed). $\mathcal{L}_\text{spec}$
with $(1+\alpha|k|)$ weighting up-weights the high-$k$ tail by an order of
magnitude at the Nyquist cutoff. Empirically (energy spectra, Appendix),
the spectral term reduces high-$k$ truncation in the prediction without
hurting low-$k$ RMSE, consistent with the theoretical prediction that
$\mathcal{L}_\text{spec}$ is a power-spectrum matcher invariant to phase
and therefore complementary to (not competing with) pixel-space
$\ell_1/\ell_2$.

\paragraph{AFNO architectural ceiling.}
ModAFNO at 12M, 127M (Slim and full) all plateau at train RMSE
$\approx 0.117$, regardless of optimiser, batch size, or epochs. We
hypothesise that the \texttt{hard\_thresholding\_fraction} mode truncation
inherent to AFNO architecturally discards high-frequency spectral content
crucial for fine-scale interpolation. A test for this hypothesis (analysing
spectral content of prediction error) is left for future work.

\input{B4_negative_results}

\input{B2_limitations_ethics}

\section{Conclusion}

We have presented a systematic architectural study of 6-hour ERA5
atmospheric time interpolation at 0.5$^\circ$ resolution ($360{\times}720$,
27 channels). Across four major architecture families (DC-AE, SwinV2,
S-DYff, ModAFNO) trained at identical 8-epoch budget on six years of ERA5,
the proposed WeatherDCAE-Skip recipe achieves $-47\%$ test RMSE versus the bicubic
baseline at only \textbf{14.4M parameters} --- outperforming $7\times$-larger
S-DYff ($-19\%$) and $3\times$-larger ModAFNO ($+49\%$). Beyond the three
known ingredients (bilinear scaffold, FiLM/AdaLN time conditioning, zero-init
U-Net skips) we introduced two novel components: (A) $\tau$-conditional
skip gates, where a small MLP on the time embedding produces per-skip
gating scalars that capture the $\tau$-dependent low-vs-high-frequency
trade-off at $<0.5$~k extra parameters; and (C) a spectral consistency
loss that linearly up-weights high wavenumbers in the FFT magnitude
error, counteracting the natural $k^{-3}$ atmospheric power-law decay
and the spectral bias of deterministic regression. Comprehensive ablations demonstrate that the bilinear scaffold is
essential (direct-mode catastrophic failure across architectures), and
a lesion study confirms that the skip pathway remains functionally
necessary at the 0.5$^\circ$+6-year scale: disabling skip gates at
inference inflates per-channel RMSE by a median of $+47\%$ on the
exact same backbone, with wind components more than doubling. The
newly proposed $\tau$-gates and spectral loss are designed to compound
on top of this strong DC-AE recipe. Code, weights, and the regridding/evaluation harness are
released for full reproducibility.

% Bibliography
\bibliographystyle{plainnat}
\bibliography{references}

\input{07_appendix}

\input{A7_multi_year_test}

\input{A8_legacy_1deg_results}

\input{B3_reproducibility}

\end{document}
